\section{总结与延伸}

\subsection{核心收获}

LangChain 这篇文章的核心论点可以归纳为一句话：\textbf{Deep Agent 的评测不能套用传统 LLM 评测的范式，需要从根本上重新设计}。具体来说：

\begin{enumerate}
  \item \textbf{测试粒度}：从``一个评测器适用所有''到``每个用例定制断言''
  \item \textbf{运行模式}：从``全量运行''到``单步 + 全轮 + 多轮''三级体系
  \item \textbf{检查维度}：从``只看输出''到``轨迹 + 响应 + 状态''三维验证
  \item \textbf{基础设施}：从``无状态评测''到``环境隔离 + API Mock''
\end{enumerate}

\subsection{与软件工程测试体系的对应关系}

Deep Agent 评测体系与经典的软件工程测试金字塔有着惊人的对应关系：

\begin{table}[H]
\centering
\caption{Deep Agent 评测与软件测试金字塔的映射}
\begin{tabular}{lll}
\toprule
\textbf{软件测试层级} & \textbf{Deep Agent 评测模式} & \textbf{建议占比} \\
\midrule
单元测试 & 单步评测（Single Step） & $\sim$50\% \\
集成测试 & 完整轮次评测（Full Turn） & $\sim$30\% \\
端到端测试 & 多轮对话评测（Multiple Turns） & $\sim$20\% \\
\bottomrule
\end{tabular}
\end{table}

\subsection{拓展阅读}

\begin{itemize}
  \item LangSmith 官方文档：\href{https://docs.smith.langchain.com/}{docs.smith.langchain.com} --- 详细的评测集成指南
  \item LangGraph 文档：\href{https://langchain-ai.github.io/langgraph/}{langchain-ai.github.io/langgraph} --- 了解 Agent 状态管理和中断机制
  \item Harbor / TerminalBench：编程 Agent 专用评测环境
  \item VCR.py：\href{https://vcrpy.readthedocs.io/}{vcrpy.readthedocs.io} --- Python HTTP 录制回放库
  \item LangChain 博客原文：\href{https://blog.langchain.com/evaluating-deep-agents-our-learnings/}{Evaluating Deep Agents: Our Learnings}
\end{itemize}

%% --- Fin del contenido --- %%

\end{document}
